% Part: incompleteness
% Chapter: theories-computability
% Section: introduction

\documentclass[../../../include/open-logic-section]{subfiles}

\begin{document}

\olfileid[hi]{inc}{tcp}{int}

\olsection{परिचय}

\begin{editorial}
इस अनुभाग को फिर से लिखा जाना चाहिए।
\end{editorial}

हमारे पास निम्न परिणाम हैं:
\begin{enumerate}
\item $\Th{Q}$ में किसी फलन के निरूपणीय होने के अर्थ की परिभाषा
  (\olref[req][int]{defn:representable-fn});
\item $\Th{Q}$ में किसी संबंध के निरूपणीय होने के अर्थ की परिभाषा
  (\olref[req][rel]{defn:representing-relations});
\item यह कहने वाला प्रमेय कि $\Th{Q}$ के निरूपणीय फलन ठीक संगणनीय फलन हैं
  (\olref[req][int]{thm:representable-iff-comp});
\item यह कहने वाला प्रमेय कि $\Th{Q}$ के निरूपणीय संबंध ठीक संगणनीय संबंध
  हैं (\olref[req][rel]{thm:representing-rels})।
\end{enumerate}

कोई {\em सिद्धांत} वाक्यों का ऐसा समुच्चय है जो निगमनात्मकतः संवृत हो;
अर्थात् उसमें यह गुण हो कि जब भी $T$,~$!A$ को सिद्ध करे, तब $!A$,~$T$ में
हो। किसी सिद्धांत को वाक्यों का ऐसा संग्रह समझना शायद सबसे अच्छा है जिसमें
इन वाक्यों से निष्पन्न होने वाली सभी बातें भी सम्मिलित हों। अब से हम
$\Th{Q}$ का उपयोग उस {\em सिद्धांत} के लिए करेंगे जिसमें
\olref[req][int]{sec} के आठ अभिगृहीतों से व्युत्पन्न वाक्यों का समुच्चय है।
याद रखें कि हम $\Th{Q}$ के सूत्रों को संख्याओं से कूटबद्ध कर सकते हैं; यदि
$!A$ ऐसा सूत्र है, तो $!A$ को कूटबद्ध करने वाली संख्या को $\Gn{!A}$ लिखें।
इस कूटलेखन के आधार पर हम अब पूछ सकते हैं कि सूत्रों के विविध समुच्चय
संगणनीय हैं या नहीं।

\end{document}
